\chapter{Base-\(K\) spectral renormalization gate}

Путь base-\(K\) сохраняет построенное четырёхмерное parent action, но
требует определить finite части determinant без обращения к наблюдаемому
vacuum scale. Проверим три естественных prescription: zeta-minimal,
геометрический scale и spectral cutoff.

\section{Общий renormalized functional}

Вдоль общей scale-direction локальная one-loop часть имеет форму
\begin{equation}
 V_{\mathrm{1loop}}(\chi)
 =
 B\chi^4
 \left(
 \log\frac{\chi^2}{\mu^2}-c
 \right),
 \qquad
 B>0,
\end{equation}
где zero-mode гейт дал
\begin{equation}
 B_0=\frac{67}{64\pi^2}
\end{equation}
до curvature-dependent determinant corrections.

Разрешённая renormalized local completion содержит
\begin{equation}
 V_{\mathrm{ren}}(\chi)
 =
 V_{\mathrm{1loop}}(\chi)
 +\lambda_2R_K\chi^2
 +\lambda_4\chi^4
 +\lambda_0R_K^2.
\end{equation}
Последний член не влияет на vacuum equation, но необходим для полной
гравитационной subtraction scheme.

\section{Zeta-scale ambiguity}

Пусть
\begin{equation}
 \mu\longmapsto e^t\mu.
\end{equation}
Тогда
\begin{equation}
 \log\frac{\chi^2}{e^{2t}\mu^2}
 =
 \log\frac{\chi^2}{\mu^2}-2t,
\end{equation}
и изменение scale точно поглощается заменой
\begin{equation}
 \lambda_4\longmapsto\lambda_4+2Bt.
\end{equation}
Следовательно, zeta determinant определяет nonlocal часть, но не отделяет
канонически \(\mu\) от finite quartic counterterm.

Даже если наложить scale-invariant условие \(\lambda_2=0\), nonzero
stationary point удовлетворяет
\begin{equation}
 \log\frac{\chi_*^2}{\mu^2}
 =
 c-\frac12-\frac{\lambda_4}{B},
\end{equation}
то есть
\begin{equation}
 \frac{\chi_*}{\mu}
 =
 \exp\left[
 \frac12
 \left(
 c-\frac12-\frac{\lambda_4}{B}
 \right)
 \right].
\end{equation}
Произвольная finite часть \(\lambda_4\) непрерывно сдвигает даже
безразмерное отношение \(\chi_*/\mu\).

\begin{nogocriterion}[Zeta-minimal no-go]
Условие «finite counterterms равны нулю» является допустимой схемой, но
не следствием spectrum. Без дополнительного normalization principle
zeta-minimal prescription не превращает \(\chi_*/\mu\) в предсказание.
\end{nogocriterion}

\section{Геометрический выбор \(\mu=R^{-1}\)}

На фиксированном unit-radius фоне можно положить
\begin{equation}
 \mu=R^{-1}.
\end{equation}
Это устраняет отдельную размерную константу и переводит vacuum equation в
условие для \(\chi R\). Однако при homothety
\begin{equation}
 R\longmapsto sR
\end{equation}
масштаб \(\mu\) также меняется. Поэтому prescription фиксирует только
dimensionless ratio и не выбирает абсолютный физический scale.

Кроме того, quartic ambiguity сохраняется:
\begin{equation}
 \chi_*R
 =
 \exp\left[
 \frac12
 \left(
 c-\frac12-\frac{\lambda_4}{B}
 \right)
 \right].
\end{equation}
Геометрическая единица длины не фиксирует finite subtraction constant.

\begin{proposition}[Geometric-scale partial pass]
Выбор \(\mu=R^{-1}\) совместим с base-\(K\) геометрией и не вводит новый
размерный параметр, но без вывода \(\lambda_4\) он не предсказывает даже
\(\chi R\). Абсолютный scale остаётся homothety-модулем.
\end{proposition}

\section{Spectral cutoff branch}

Обычный spectral action вводит cutoff \(\Lambda\) и moments spectral
function. Тогда bare coefficients имеют вид
\begin{equation}
 \lambda_2=\lambda_2(f_2,\Lambda),
 \qquad
 \lambda_4=\lambda_4(f_0).
\end{equation}
Эта схема действительно может задать renormalization boundary condition,
но \(f_0,f_2\) и \(\Lambda\) являются новыми непрерывными входами, если не
выведены из динамики dilaton/radion sector. Нормировка scalar kinetic term
устраняет один общий множитель, но не фиксирует независимо все relevant
moments.

\begin{nogocriterion}[Spectral-moment input rule]
Использование \(\Lambda,f_0,f_2\) для выбора vacuum допустимо только как
расширенная модель с явно объявленными train inputs. Оно не является
parameter-free завершением текущей версии III.
\end{nogocriterion}

\section{Вердикт subtraction gate}

\begin{center}
\begin{tabular}{p{0.25\textwidth}p{0.27\textwidth}p{0.32\textwidth}}
\toprule
Prescription & Что фиксирует & Оставшаяся свобода \\
\midrule
zeta-minimal & конкретную вычислительную схему & \(\mu\) и эквивалентная finite \(\lambda_4\) \\
геометрическая \(\mu=R^{-1}\) & отсутствие нового dimensionful input & finite \(\lambda_4\), homothety scale \\
spectral cutoff & bare boundary condition & \(\Lambda,f_0,f_2\) как continuous inputs \\
\bottomrule
\end{tabular}
\end{center}

\begin{theorem}[Base-\(K\) renormalization no-go]
Ни zeta-minimal, ни геометрический scale, ни обычный spectral cutoff не
фиксируют parameter-free vacuum текущего parent action. Положительный
\(B_0\) обеспечивает возможность Coleman--Weinberg минимума после выбора
scheme, но его положение остаётся renormalization datum.
\end{theorem}

Следовательно, base-\(K\) путь не закрывается как абсолютная
предсказательная теория. Его корректный статус --- renormalizable
four-dimensional effective parent action с выведенными относительными
коэффициентами и одним открытым scale-setting condition.

\begin{protocol}[Следующий конструктивный выбор]
Остаются два непересекающихся продолжения:
\begin{enumerate}
 \item принять один scale-setting observable как явный train input и
 проверять только dimensionless blind ratios;
 \item искать новую динамическую symmetry, RG fixed point или anomaly
 condition, которая фиксирует finite \(\lambda_4/B\) без mass data.
\end{enumerate}
До такого выбора запрещено объявлять абсолютный \(\chi_*\) предсказанием.
\end{protocol}

Машинный аудит сохранён в
\path{s2t_v3_base_k_spectral_renormalization_results.json}.